% !TeX program = XeLaTeX
% !TeX root = ../pujavidhanam.tex
\chapt{श्री-ऋषिपञ्चमी-व्रतम्}

\centerline{\small{(मूलम्—श्रीव्रतरत्नाकरः, प्रथमो भागः; उद्यापनविधिः ब्रह्माण्डपुराणोक्तः)}}

\input{purvanga/vighneshwara-puja}

\sect{प्रधान-पूजा — सप्तर्षि-अरुन्धती-पूजा}

\twolineshloka*
{शुक्लाम्बरधरं विष्णुं शशिवर्णं चतुर्भुजम्}
{प्रसन्नवदनं ध्यायेत् सर्वविघ्नोपशान्तये}

भाद्रपद-शुक्ल-पञ्चम्यां मध्याह्नव्यापिन्यां कर्तव्यमिदं व्रतम्। प्रातः स्नानानन्तरं
अपामार्ग-समिधा (अथवा सप्तसंख्याभिः) दन्तधावनं कुर्यात्।

\dnsub{दन्तधावन-प्रार्थना}
\twolineshloka*
{आयुर्बलं यशो वर्चः प्रजाः पशुवसूनि च}
{ब्रह्म प्रज्ञां च मेधां च त्वन्नो देहि वनस्पते}

\twolineshloka*
{मुखदुर्गन्धनाशाय दन्तानां च विशुद्धये}
{ष्ठीवनाय च गात्राणां कुर्वेऽहं दन्तधावनम्}

अनेन दन्तान् संशोध्य मृत्स्नापूर्वकं स्नायात्। रजस्वला-दोष-निवृत्त्यर्थं रक्तवस्त्रदानं
कुर्यात्। ततः पञ्चगव्य-सम्मेलनं कृत्वा प्रतिमां शोधयेत्।

\dnsub{पञ्चगव्य-मन्त्राः}
गायत्र्या गोमूत्रम्, "गन्धद्वारा"मिति गोमयम्, "आप्यायस्वे"ति क्षीरम्,
"दधिक्राव्णो" इति दधि, "शुक्रमसी"त्याज्यम्, "देवस्य त्वे"ति कुशोदकम् — एभिः
दर्भैः प्रणवेनालोड्य, प्रणवेनाभिमन्त्र्य पञ्चगव्येन प्रतिमाशुद्धिं कृत्वा प्राशयेत्।

\twolineshloka*
{यत्त्वगस्थिगतं पापं देहे तिष्ठति मामके}
{प्राशनं पञ्चगव्यस्य दहत्वग्निरिवेन्धनम्}

गृहं गत्वा गोमयेनोपलिप्तां रङ्गवल्लीशोभितां वेदिं सर्वतोभद्रमण्डलयुक्तां
कुर्यात्, अव्रणं सजलं ताम्रं वा मृण्मयं कुम्भं वस्त्रसंयुक्तं
पञ्चरत्नफलगन्धाक्षतैर्युतं सहिरण्यं संस्थाप्य, कण्ठे वेष्टितं वस्त्रेणाच्छाद्य,
अष्टदलं तत्र लिखित्वा तस्मिन् सप्त ऋषीन् दिव्यान् भक्तियुक्तः प्रपूजयेत्।

\dnsub{सङ्कल्पः}

ममोपात्त-समस्त-दुरित-क्षयद्वारा श्री-परमेश्वर-प्रीत्यर्थं शुभे शोभने मुहूर्ते अद्य ब्रह्मणः
द्वितीयपरार्धे श्वेतवराहकल्पे वैवस्वतमन्वन्तरे अष्टाविंशतितमे कलियुगे प्रथमे पादे
जम्बूद्वीपे भारतवर्षे भरतखण्डे मेरोः दक्षिणे पार्श्वे शकाब्दे अस्मिन् वर्तमाने व्यावहारिकाणां
प्रभवादीनां षष्ट्याः संवत्सराणां मध्ये \blank\see{app:samvatsara_names} नाम संवत्सरे दक्षिणायने
वर्ष-ऋतौ भाद्रपद-मासे शुक्लपक्षे पञ्चम्यां शुभतिथौ \blank{} वासरयुक्तायां मध्याह्नव्यापिन्यां
\blank\see{app:nakshatra_names} नक्षत्र \blank\see{app:yoga_names}-योग \blank\see{app:karanam_names}-करण-युक्तायां च एवं गुण विशेषण विशिष्टायाम्
अस्याम् पञ्चम्यां शुभतिथौ
मम इहजन्मनि जन्मान्तरे च स्त्री-जन्म-प्रादुर्भाव-सम्भूत-रजोदर्शनादारभ्य रजोनिवृत्तिपर्यन्तं
मध्यवर्तिनि काले ऋतुमत्यां सत्यां ज्ञानतः अज्ञानतश्च आचरित-सकल-प्रतिबन्धक-दोष-निवृत्त्यर्थं
मम इहजन्मनि पूर्वजन्मनि जन्मान्तरे च सम्पादितानां ज्ञानाज्ञानकृतमहा\-पातकचतुष्टय-व्यतिरिक्तानां
रहस्यकृतानां प्रकाशकृतानां सर्वेषां पापानां सद्य अपनोदनद्वारा सकल-पापक्षयार्थं
सौभाग्य-सन्तान-सौख्य-आयुरारोग्य-सिद्ध्यर्थम् अरुन्धती-सहित-कश्यपादि-सप्तर्षि-प्रसादसिद्ध्यर्थं
यथाशक्ति-ध्यानावाहनादिषोडशोपचारैः अरुन्धती-सहित-कश्यपादि-सप्तर्षि-पूजां करिष्ये। तदङ्गं कलशपूजां च करिष्ये।

श्रीविघ्नेश्वराय नमः यथास्थानं प्रतिष्ठापयामि। शोभनार्थे क्षेमाय पुनरागमनाय च।\\
(गणपति-प्रसादं शिरसा गृहीत्वा)

\input{purvanga/aasana-puja}

\input{purvanga/ghanta-puja}

\begingroup
\renewcommand{\devaName}{ऋषि}
\input{purvanga/kalasha-puja}
\endgroup

\input{purvanga/aatma-puja}

\input{purvanga/pitha-puja}

\input{purvanga/guru-dhyanam}

\sect{प्रत्येक-ऋषि-पूजा}

व्रतारम्भे आचार्य-ऋत्विग्वरणं कृत्वा, कुम्भेषु वरुणमावाह्य, तदुपरि प्रतिमासु
कश्यपादीन् आवाहयेत्। आदौ प्रत्येकम् एकैकं ऋषिं सम्पूज्य, ततः समया सर्वान् ऋषीन् पूजयेत्।

\dnsub{१. कश्यप-पूजा}
\twolineshloka*
{जटिलं भस्मदिग्धाङ्गं साक्षसूत्रकमण्डलुम्}
{वल्कलाजिनकौपीनं सातपत्रं सदण्डकम्}
\textbf{अदिति-सहितं कश्यपं ध्यायामि।}

ऋषेर्मन्त्रकृतामित्यस्य मन्त्रस्य कश्यप ऋषिः। कश्यपो देवता। पङ्क्तिश्छन्दः।
कश्यपावाहने विनियोगः।
\twolineshloka*
{ऋषेर्मन्त्रकृतां स्तोमैः कश्यपोद्वर्धयन् गिरः}
{सोमं नमस्य राजानं यो जज्ञे वीरुधां पतिः}
\textbf{अदिति-सहितं कश्यपम् आवाहयामि।}

कश्यपाय आसनम्। मरीचिपुत्राय पाद्यम्। सूर्यजनकाय अर्घ्यम्। इन्द्रताताय
आचमनीयम्। वामनगुरवे मधुपर्कम्। अदितिभर्त्रे पञ्चामृतस्नानम्। ब्रह्मनेत्राय
शुद्धोदकस्नानम्। देवप्रियाय वस्त्रद्वयम्। सुवर्चसे यज्ञोपवीतम्। विशालाक्षाय
आभरणानि। सर्वविदे गन्धम्। महते अक्षतान् — इति क्रमेण समर्पयामि।

कश्यपाय नमः; वामनगुरवे नमः; विशालाक्षाय नमः; मरीचिपुत्राय नमः; अदितिभर्त्रे नमः; सर्वविदे नमः; कविजनकाय नमः; देवप्रियाय नमः; महते नमः; इन्द्रताताय नमः; सुवर्चसे नमः; विमलाय नमः।
\textbf{इति पुष्पाणि समर्पयामि।} धीमते धूपम्। अकल्मषाय दीपम्। सुमनोरथाय
नैवेद्यम्। जटिने ताम्बूलम्। दीप्ततेजसे नीराजनम्। दण्डिने मन्त्रपुष्पाणि समर्पयामि।

\dnsub{२. अत्रि-पूजा}
\twolineshloka*
{मरीचितनयं देवं साक्षसूत्रकमण्डलुम्}
{अनसूयापतिं शान्तम् अत्रिं निष्कल्मषं मुनिम्}
\textbf{अनसूया-समेतम् अत्रिं ध्यायामि।}

स्वर्भानोरित्यस्य मन्त्रस्य मरीचिपुत्रोऽत्रिः ऋषिः। अत्रिर्देवता। त्रिष्टुप् छन्दः।
अत्र्यावाहने विनियोगः।
\twolineshloka*
{स्वर्भानो रथ यदिन्द्रमाया आवोदिवा वर्तमाना आवाहन्}
{गूळ्हं सूर्यं तमसापवृतेन तुरीयेण ब्रह्मणा विन्ददत्रिः}
\textbf{अनसूया-समेतम् अत्रिम् आवाहयामि।}

अत्रये आसनम्। धीमते पाद्यम्। योगीन्द्राय अर्घ्यम्। दत्तात्रेयप्रियाय आचमनीयम्।
चन्द्रताताय मधुपर्कम्। अनसूयापतये पञ्चामृतम्। नित्यशुद्धाय शुद्धोदकम्।
वल्कलधारिणे वस्त्रम्। ब्रह्मिष्ठाय यज्ञोपवीतम्। दीप्ताय आभरणम्। दण्डधारिणे
गन्धम्। त्रिमूर्तिजनकाय अक्षतान् — इति क्रमेण समर्पयामि।

अत्रये नमः; धीमते नमः; योगीन्द्राय नमः; दत्तात्रेयप्रियाय नमः; त्रिमूर्तिजनकाय नमः; चन्द्रताताय नमः; द्युतिमते नमः; मरीचिपुत्राय नमः।
\textbf{इति पुष्पाणि समर्पयामि।} सुगन्धदेहाय धूपम्। द्युतिमते दीपम्।
नित्यतृप्ताय नैवेद्यम्। विरागिणे ताम्बूलम्। अर्चिष्मते नीराजनम्। भ्राजिष्णवे
मन्त्रपुष्पम्। संयमिने प्रदक्षिणम्। दमिने नमस्कारान्। कल्मषघ्नाय पुनरर्घ्यं
समर्पयामि। अत्रये सर्वोपचारान् समर्पयामि।

\dnsub{३. भरद्वाज-पूजा}
\twolineshloka*
{भरद्वाजं महाशान्तं सुशीलापतिम् ऊर्जितम्}
{अक्षस्रग्गन्धहस्तं च मुनिम् आङ्गिरसं भजे}
\textbf{सुशीला-सहितं भरद्वाजं ध्यायामि।}

अविन्दन्नित्यस्य मन्त्रस्य सौज्जिको भारद्वाज ऋषिः। भरद्वाजो देवता। त्रिष्टुप्छन्दः।
भरद्वाजावाहने विनियोगः।
\twolineshloka*
{अविन्दतेति हि हितं यदासीद्यज्ञस्य धाम परमं गुहा यत्}
{धातुर्द्युतानात्सवितुश्च विष्णोर्भरद्वाजो बृहदाचक्रे अग्नेः}
\textbf{सुशीला-सहितं भरद्वाजम् आवाहयामि।}

भरद्वाजाय आसनम्। भ्राजिष्णवे पाद्यम्। विदुषे अर्घ्यम्। अकल्मषाय आचमनीयम्।
वल्कलधारिणे मधुपर्कम्। पवित्रदेहाय पञ्चामृतस्नानम्। परिशुद्धाय
शुद्धोदकस्नानम्। ज्ञानिने वस्त्रम्। मुनये यज्ञोपवीतम्। ब्रह्मचारिणे आभरणानि।
त्रिकालज्ञानिने गन्धम्। मुनये अक्षतान् — इति क्रमेण समर्पयामि।

भरद्वाजाय नमः; वल्कलधारिणे नमः; मुनये नमः; भ्राजिष्णवे नमः; पवित्रदेहाय नमः; ब्रह्मचारिणे नमः; विदुषे नमः; परिशुद्धाय नमः; त्रिकालज्ञाय नमः; अकल्मषाय नमः; ज्ञानिने नमः; सुप्रसिद्धाय नमः।
\textbf{इति पुष्पाणि समर्पयामि।} सुनासिकाय धूपम्। ज्ञानिने दीपम्।
अवाप्तसकलकामाय नैवेद्यम्। सुमनस्काय ताम्बूलम्। तेजसे नीराजनम्। भ्राजिष्णवे
मन्त्रपुष्पम्। जितक्रोधाय आत्मप्रदक्षिणम्। अवन्दिने नमस्कारान्। मुमुक्षवे
पुनरर्घ्यम्। भरद्वाजाय सर्वोपचारान् समर्पयामि।

\dnsub{४. विश्वामित्र-पूजा}
\twolineshloka*
{हिरण्यगर्भसम्भूतं जटामकुटधारिणम्}
{कुमुद्वतीपतिं शान्तं विश्वामित्रमहं भजे}
\textbf{कुमुद्वती-सहितं विश्वामित्रं ध्यायामि।}

इमे भोजा इत्यस्य मन्त्रस्य प्रगाथो विश्वामित्रः ऋषिः। त्रिष्टुप्छन्दः। गाधिपुत्रो
विश्वामित्रो देवता। विश्वामित्रावाहने विनियोगः।
\twolineshloka*
{इमे भोजा अङ्गिरसो विरूपा दिवस्पुत्रासो असुरस्य वीराः}
{विश्वामित्राय ददतो मघानि सप्त प्रतिरन्तो न आयुः}
\textbf{कुमुद्वती-सहितं विश्वामित्रम् आवाहयामि।}

विश्वामित्राय आसनम्। दग्धकल्मषाय पाद्यम्। शुद्धबुद्धये अर्घ्यम्। उग्रतपसे
आचमनीयम्। अव्यग्राय मधुपर्कम्। दण्डपाणये पञ्चामृतम्। कमण्डलुधराय
शुद्धोदकम्। अजिनाम्बराय वस्त्रम्। तेजस्विने यज्ञोपवीतम्। रम्याय आभरणानि।
कुशिकपौत्राय गन्धम्। वरदाय अक्षतान् — इति क्रमेण समर्पयामि।

मुमुक्षवे नमः; गायत्रीप्रियाय नमः; वेदशास्त्रार्थतत्त्वज्ञाय नमः; वीराय नमः; अमायिने नमः; ब्रह्मर्षये नमः; दग्धकल्मषाय नमः; अनघतेजसे नमः; बुद्धिप्रदाय नमः; उग्रतपसे नमः; अव्यग्राय नमः; दण्डपाणये नमः; कमण्डलुधराय नमः; अजिनाम्बराय नमः; तेजस्विने नमः; रम्याय नमः; कुशिकपौत्राय नमः; वरदाय नमः; गाधितनूभवाय नमः।
\textbf{इति पुष्पाणि समर्पयामि।} कमण्डलुधराय धूपम्। तेजसे दीपम्। रम्याय
नैवेद्यम्। कौशिकाय ताम्बूलम्। वरदाय नीराजनम्। गाधितनयाय मन्त्रपुष्पम्।
रुद्रप्रियाय प्रदक्षिणम्। गायत्रीजपानन्दाय नमस्कारान्। वेदशास्त्रार्थतत्त्वज्ञाय
पुनरर्घ्यम्। विश्वामित्राय सर्वोपचारान् समर्पयामि।

\dnsub{५. गौतम-पूजा}
\twolineshloka*
{रहूगणात्मजं शान्तं त्रिपुण्ड्राङ्कितमस्तकम्}
{अक्षस्रक्कुण्डले दण्डान् दधानं गौतमं भजे}
\textbf{अहल्या-सहितं गौतमं ध्यायामि।}

अभित्वा गौतम इत्यस्य मन्त्रस्य गौतमः ऋषिः। गायत्रीछन्दः। गौतमो देवता।
गौतमावाहने विनियोगः।
\twolineshloka*
{अभि त्वा गौतमा गिरा जातवेदो विचर्षणे}
{द्युम्नैरभि प्र णोनुमः}
\textbf{अहल्या-पतिं गौतमम् आवाहयामि।}

गौतमाय आसनम्। भस्मोद्धूलितविग्रहाय पाद्यम्। रुद्राक्षमालिने अर्घ्यम्।
रुद्रप्रियाय आचमनीयम्। देवपूजानिरताय मधुपर्कम्। अहल्यापतये पञ्चामृतम्।
ध्वस्तपापाय शुद्धोदकस्नानम्। कमण्डलुधराय वस्त्रम्। सुवर्चसे यज्ञोपवीतम्।
निष्कल्मषाय आभरणानि। निरुपद्रवाय गन्धम्। निरहङ्काराय अक्षतान् — इति
क्रमेण समर्पयामि।

गौतमाय नमः; ध्वस्तपापाय नमः; भस्मोद्धूलितविग्रहाय नमः; कमण्डलुधराय नमः; निरहङ्काराय नमः; वल्कलाम्बराय नमः; नित्यतृप्ताय नमः; शिवलिङ्गार्चकाय नमः; निष्कल्मषाय नमः; निर्मोहाय नमः; निरुपद्रवाय नमः; रुद्राक्षमालिने नमः; सुवर्चसे नमः; अहल्यापतये नमः।
\textbf{इति पुष्पाणि समर्पयामि।} विमलाय धूपम्। द्युतिमते दीपम्। नित्यतृप्ताय
नैवेद्यम्। अक्लेशाय ताम्बूलम्। वरदाय नीराजनम्। जटिने मन्त्रपुष्पम्। संयमिने
प्रदक्षिणम्। मुमुक्षवे नमस्कारान्। गौतमाय पुनरर्घ्यं सर्वोपचारान् च समर्पयामि।

\dnsub{६. जमदग्नि-पूजा}
\twolineshloka*
{शान्तं जितारिषड्वर्गं भृगुपुत्रं महाद्युतिम्}
{दण्डाक्षसूत्रपाणिं च जमदग्निं नमाम्यहम्}
\textbf{रेणुका-सहितं जमदग्निं ध्यायामि।}

प्रसूतो भक्षमित्यस्य मन्त्रस्य प्रगाथो जमदग्निः ऋषिः। जगतीछन्दः। जमदग्निर्देवता।
जमदग्न्यावाहने विनियोगः।
\twolineshloka*
{प्रसूतो भक्षं करामः स्तोमं चेमं प्रथमः सूरिरुन्मृजे}
{सुते सातेन यद्यगामं वां प्रति विश्वामित्र जमदग्नीदमे}
\textbf{रेणुकापतिं जमदग्निम् आवाहयामि।}

जमदग्नये आसनम्। निर्मदाय पाद्यम्। वीतरागाय अर्घ्यम्। कान्तिमते आचमनीयम्।
शान्ताय मधुपर्कम्। दान्ताय पञ्चामृतम्। शुचये शुद्धोदकम्। रामप्रियाय वस्त्रम्।
अकामाय यज्ञोपवीतम्। भक्तप्रियाय आभरणानि। दण्डिने गन्धम्। चण्डाय अक्षतान्
— इति क्रमेण समर्पयामि।

जमदग्नये नमः; कमण्डलुधराय नमः; इष्टाय नमः; साक्षस्रजे नमः; विमलाय नमः; पुष्टाय नमः; दयालवे नमः; कल्मषापहराय नमः; तुष्टाय नमः; दोषवर्जिताय नमः; इष्टप्रदाय नमः; त्वगम्बराय नमः।
\textbf{इति पुष्पाणि समर्पयामि।} दोषवर्जिताय धूपम्। ज्ञानिने दीपम्। विमलाय
नैवेद्यम्। कल्मषापहराय ताम्बूलम्। इष्टप्रदाय नीराजनम्। शिष्टाय मन्त्रपुष्पम्।
श्रेष्ठाय प्रदक्षिणम्। पुष्टाय नमस्कारान्। विशिष्टाय पुनरर्घ्यम्। जमदग्नये
सर्वोपचारपूजां समर्पयामि।

\dnsub{७. वसिष्ठ-पूजा}
\twolineshloka*
{तपोनिधिं दयासिन्धुं ब्रह्मण्यं ब्रह्मसम्भवम्}
{स्रग्गन्धपाणिं वरदं वसिष्ठं प्रणमाम्यहम्}
\textbf{अरुन्धती-सहितं वसिष्ठं ध्यायामि।}

श्वित्यञ्चोमा इत्यस्य मन्त्रस्य मैत्रावरुणो वसिष्ठः ऋषिः। त्रिष्टुप्छन्दः। वसिष्ठो
देवता। वसिष्ठावाहने विनियोगः।
\twolineshloka*
{श्वित्यञ्चो मा दक्षिणतस्कपर्दा धियं जिन्वासो अभि हि प्र मन्दुः}
{उत्तिष्ठन्वो वृषणो युध्यमाना नमेदूरादवितवे वसिष्ठाः}
\textbf{अरुन्धतीपतिं वसिष्ठम् आवाहयामि।}

वसिष्ठाय आसनम्। वरिष्ठाय पाद्यम्। गरिष्ठाय अर्घ्यम्। धर्मिष्ठाय आचमनीयम्।
अरुन्धतीपतये मधुपर्कम्। दीप्तिमते पञ्चामृतस्नानम्। त्रिकालज्ञाय
शुद्धोदकस्नानम्। योगिने वस्त्रम्। रामप्रियाय यज्ञोपवीतम्।
शरणागतवत्सलाय आभरणम्। मनोरथफलप्रदाय अक्षतान् — इति क्रमेण समर्पयामि।

वसिष्ठाय नमः; अरुन्धतीपतये नमः; रामप्रियाय नमः; वरिष्ठाय नमः; दीप्तिमते नमः; शरणागतवत्सलाय नमः; गरिष्ठाय नमः; त्रिकालज्ञाय नमः; मनोरथफलप्रदाय नमः; धर्मिष्ठाय नमः; योगिने नमः; भक्ताभीष्टप्रदाय नमः।
\textbf{इति पुष्पाणि समर्पयामि।} निष्कल्मषाय धूपम्। मैत्रावरुणाय दीपम्।
ब्रह्मचारिणे नैवेद्यम्। विरागिणे ताम्बूलम्। ज्ञानिने नीराजनम्। जटिने पुनरर्घ्यम्।
वेदविदे मन्त्रपुष्पम्। दान्ताय प्रदक्षिणम्। शान्ताय नमस्कारान्। ब्रह्मिष्ठाय
पुनरर्घ्यम्। वसिष्ठाय सर्वोपचारपूजाः समर्पयामि।

\dnsub{८. अगस्त्य-पूजा}
\twolineshloka*
{लोपामुद्रापतिं देवम् अगस्त्यम् ऋषिसत्तमम्}
{दक्षिणाशापतिं रामतारकप्रियम् ऊर्जितम्}
\textbf{लोपामुद्रा-समेतम् अगस्त्यं ध्यायामि।}

अगस्त्यः खनमान इत्यस्य मन्त्रस्य अगस्त्यः ऋषिः। त्रिष्टुप्छन्दः। अगस्त्यो देवता।
अगस्त्यावाहने विनियोगः।
\twolineshloka*
{अगस्त्यः खनमानः खनित्रैरुभे वर्णावृषिरुग्रः पुपोष}
{सत्या देवेष्वाशिषो जगाम स्वादित्संगृभ्ण उभयाँ अमेने}
\textbf{लोपामुद्रा-समेतम् अगस्त्यम् आवाहयामि।}

लोपामुद्रापतये आसनम्। आत्मवते अर्घ्यम्। इच्छापूर्तये पाद्यम्।
ईषणात्रयवर्जिताय आचमनीयम्। उत्तमाय मधुपर्कम्। ऊर्जिततपसे
पञ्चामृतस्नानम्। ऋग्यजुस्सामपारगाय शुद्धोदकस्नानम्। ऋणमोचनाय वस्त्रम्।
लुप्तपापाय यज्ञोपवीतम्। ओङ्कारप्रियाय आभरणानि। उदाराय अक्षतान् — इति
क्रमेण समर्पयामि।

अगस्त्याय नमः; साक्षिणे नमः; लोकपूजिताय नमः; सनातनाय नमः; ताराय नमः; तापसप्रियाय नमः; लोपामुद्रापतये नमः; रामप्रियाय नमः; विन्ध्यगुरवे नमः; कुम्भसम्भवाय नमः; रञ्जिततपसे नमः; वीतरागाय नमः; समुद्रपायिने नमः; तारकाय नमः; विष्णुरताय नमः; विश्वात्मने नमः।
\textbf{इति पुष्पाणि समर्पयामि।} गन्धज्ञाय धूपम्। ज्योतिष्मते दीपं दर्शयामि।
सारग्राहिणे नैवेद्यम्। शुचये ताम्बूलम्। प्रकाशरूपाय नीराजनम्। धनदाय दक्षिणाम्।
सुवर्चसे मन्त्रपुष्पम्। सुरेन्द्रवन्दिताय प्रदक्षिणम्। ऋषिवर्याय नमस्कारान्।
प्रसन्नाय पुनरर्घ्यम्। श्री-अगस्त्याय सर्वोपचारपूजाः समर्पयामि।

\emph{(अगस्त्यः ऋषिपञ्चमीव्रतस्य मूल-सप्तर्षि-सप्तकयोः अन्तर्गतो न भवति; अत्र
आतिथ्येन पूज्यते।)}

\dnsub{९. अरुन्धती-पूजा}
अरुन्धत्यै ध्यानम्। वसिष्ठपत्न्यै आवाहनम्। सत्यै आसनम्। पतिव्रतायै पाद्यम्।
सौमङ्गल्यहेतवे अर्घ्यम्। तपस्विन्यै आचमनीयम्। मुनिधर्मप्रवर्तिन्यै मधुपर्कम्।
नित्यनिर्मलायै स्नानम्। मानिन्यै वस्त्रम्। गोप्यै कञ्चुकम्। धर्मरतायै गन्धम्।
अक्षतव्रतायै अक्षतान्। सुमङ्गल्यै हरिद्राचूर्णम्। वसिष्ठप्रियायै कुङ्कुमम्।
सुगुणोज्ज्वलायै आभरणानि। पावन्यै पुष्पाणि — इति क्रमेण समर्पयामि।
हविर्गन्धप्रियायै धूपम्। ज्योतिर्मयायै दीपम्। धर्मब्रह्मानन्दवत्यै नैवेद्यम्।
सर्वभूतानुरागिण्यै ताम्बूलम्। ऋतुदक्षिणायै दक्षिणाम्। ऋषिपत्नीसमाश्रितायै
नीराजनम्। मन्त्रज्ञायै मन्त्रपुष्पम्। शिंशुमारगत्यै प्रदक्षिणम्।
सर्वलोकनमस्कृतायै नमस्कारान्। वसिष्ठकुटुम्बिन्यै पुनरर्घ्यं समर्पयामि।

\twolineshloka*
{स्मृतिमात्रेण भक्तानां सर्वपापविनोदिनि}
{सूक्ष्मात्सूक्ष्मतरे मातर्दर्शनात्पुण्यदायिनि}

\twolineshloka*
{पतिव्रताग्रगण्ये त्वं सौमङ्गल्यसुखप्रदे}
{अरुन्धति महादेवि ब्रह्मशक्ते चिरन्तनि}

\twolineshloka*
{जय त्वं नित्यकल्याणि वसिष्ठैककुटुम्बिनि}
{पावनत्वं च सौभाग्यं सुखं शाश्वतमेव च}

\onelineshloka*
{महादेवि जगद्वन्द्ये व्रतसम्पूर्णतां कुरु}

\sect{सामूहिक सप्तर्षि-अरुन्धती-पूजा}

\twolineshloka*
{आगच्छन्तु महाभागाश्चतुर्वेदपरायणाः}
{यावद्व्रतमिदं कुर्वे कृपया भवतामहम्}
\textbf{अरुन्धती-सहित-कश्यपादि-सप्तर्षीन् आवाहयामि।} प्राणप्रतिष्ठां कुर्यात्।

\twolineshloka*
{ऋग्यजुस्सामवेदानां स्वरूपेभ्यो नमो नमः}
{पुराणपुरुषेभ्यो हि देवर्षिभ्यो नमो नमः}
\textbf{अरुन्धतीसहित-कश्यपादि-सप्तर्षिभ्यो नमः आसनं समर्पयामि।}

\twolineshloka*
{गन्धपुष्पाक्षतैर्युक्तं पाद्यं गृह्णन्तु वै द्विजाः}
{प्रसादं कुर्वतां प्रीतास्तुष्टास्सन्तु सदा मम}
\textbf{...भ्यो नमः पाद्यं समर्पयामि।}

\twolineshloka*
{नभस्ये शुक्लपञ्चम्यामर्चिता ऋषिसत्तमाः}
{दहन्तु पापं मे सर्वे गृह्णन्त्वर्घ्यं नमो नमः}
\textbf{...भ्यो नमः अर्घ्यं समर्पयामि।}

\twolineshloka*
{लोकानां तुष्टिकर्तारो यूयं सर्वे तपोधनाः}
{नमो वो धर्मविज्ञेभ्यो महर्षिभ्यो नमो नमः}
\textbf{...भ्यो नमः आचमनीयं समर्पयामि।}

\twolineshloka*
{पयो दधि घृतं चैव शर्करामधुसंयुतम्}
{पञ्चामृतेन स्नपनं करोमि मुनिसत्तमाः}
\textbf{...भ्यो नमः पञ्चामृतस्नानं समर्पयामि।}

\twolineshloka*
{मन्दाकिन्यास्समानीतं यमुनागौतमीजलम्}
{कृष्णा-नर्मदा-गोदाभ्यः स्नानार्थमाहृतम्}
\textbf{...भ्यो नमः शुद्धोदकस्नानं समर्पयामि। स्नानानन्तरम् आचमनीयं समर्पयामि।}

\twolineshloka*
{सर्वे नित्यं तपोनिष्ठा ब्रह्मज्ञास्सत्यवादिनः}
{वस्त्राणि प्रतिगृह्णन्तु मुक्तिदास्सन्तु मे सदा}
\textbf{...भ्यो नमः वस्त्रं समर्पयामि।}

\twolineshloka*
{नानामन्त्रैस्समुद्भूतं त्रिवृतं ब्रह्मसूत्रकम्}
{प्रत्येकं च प्रयच्छामि मुनयः प्रतिगृह्यताम्}
\textbf{...भ्यो नमः उपवीतानि समर्पयामि।}

\twolineshloka*
{कुङ्कुमागरुकर्पूरसुगन्धैर्मिश्रितं शुभम्}
{गन्धाढ्यं चन्दनं दिव्यं गृह्णन्तु मुनिसत्तमाः}
\textbf{...भ्यो नमः गन्धं समर्पयामि।}

\twolineshloka*
{शुभ्राक्षताश्च सम्पूर्णाः प्रक्षाल्य च नियोजिताः}
{शोभायै वो मया दत्ता गृह्यन्तां मुनिसत्तमाः}
\textbf{...भ्यो नमः अक्षतान् समर्पयामि।}

\twolineshloka*
{मालतीचम्पकामल्लीकेतक्यादीनि वै द्विजाः}
{समर्पयामि पुष्पाणि पूजार्थं वः पदाम्बुजे}
\textbf{...भ्यो नमः पुष्पमालां समर्पयामि।}

\dnsub{अङ्ग-पूजा}
\begin{longtable}{ll@{~नमः — }l@{~पूजयामि}}
    १. & तीर्थपादेभ्यो & पादौ\\
    २. & सर्वभूतहितेभ्यो & गुल्फौ\\
    ३. & विमलचरित्रेभ्यो & जङ्घे\\
    ४. & सर्वलोकहितेभ्यो & ऊरू\\
    ५. & अजिनधारिभ्यो & कटिं\\
    ६. & जितेन्द्रियेभ्यो & गुह्यम्\\
    ७. & योगीश्वरेभ्यो & नाभिम्\\
    ८. & परार्थैकप्रयोजनेभ्यो & पार्श्वे\\
    ९. & सर्वशास्त्रविशारदेभ्यो & उदरम्\\
    १०. & परोपकारिभ्यो & हृदयम्\\
    ११. & रुद्राक्षधारिभ्यो & कण्ठम्\\
    १२. & सर्वामरवन्दितेभ्यो & नासिकाम्\\
    १३. & श्रुतिपारगेभ्यो & श्रोत्रे\\
    १४. & तत्त्वविद्याविशारदेभ्यो & ललाटम्\\
    १५. & जटामण्डलमण्डितेभ्यो & शिरः\\
    १६. & सप्तर्षिभ्यो & सर्वाण्यङ्गानि\\
\end{longtable}

\dnsub{ऋषि-अष्टोत्तरशत-नामावलिः}
\centerline{(प्रतिनाम "नमः" इति अनुसन्धेयम्)}

\begin{multicols}{3}
\begin{enumerate}
\item ब्रह्मर्षिभ्यो
\item अध्वर्युभ्यो
\item वेदविद्भ्यो
\item यज्वभ्यो
\item तपस्विभ्यो
\item यज्ञदीक्षितेभ्यो
\item स्वरूपसुखिभ्यो
\item प्रवृत्तिधर्मपालकेभ्यो
\item निवृत्तिधर्मदर्शकेभ्यो
\item महात्मभ्यो
\item पूतेभ्यो
\item भगवत्प्रसादिभ्यो
\item मान्येभ्यो
\item पुरातनेभ्यो
\item देवगुरुभ्यो
\item ब्रह्मचर्यरतेभ्यो
\item सृष्टिकर्तृभ्यो
\item लोकगुरुभ्यो
\item सिद्धेभ्यो
\item स्थितिकर्तृभ्यो
\item सर्ववन्द्येभ्यो
\item कर्मठेभ्यो
\item लयकर्तृभ्यो
\item सर्वपूज्येभ्यो
\item योगिभ्यो
\item जपकर्तृभ्यो
\item गृहिभ्यो
\item अग्निहोत्रपरायणेभ्यो
\item होतृभ्यो
\item सूत्रकृद्भ्यो
\item सत्यव्रतेभ्यो
\item प्रस्तोतृभ्यो
\item भाष्यकृद्भ्यो
\item धर्मात्मभ्यो
\item प्रतिहर्तृभ्यो
\item महिमसिद्धेभ्यो
\item उद्गातृभ्यो
\item ज्ञानसिद्धेभ्यो
\item ब्रह्मण्येभ्यो
\item धर्मप्रवर्तकेभ्यो
\item निर्दुष्टेभ्यो
\item ब्रह्मास्त्रविद्भ्यो
\item आचारप्रवर्तकेभ्यो
\item शमधनेभ्यो
\item ब्रह्मदण्डधरेभ्यो
\item सम्प्रदायप्रवर्तकेभ्यो
\item तपोधनेभ्यो
\item ब्रह्मशीर्षविद्भ्यो
\item अनुशासितृभ्यो
\item शापशक्तेभ्यो
\item गायत्रीसिद्धेभ्यो
\item वेदवेदान्तपारगेभ्यो
\item मन्त्रमूर्तिभ्यो
\item सावित्रीसिद्धेभ्यो
\item वेदाङ्गप्रचारकेभ्यो
\item अष्टाङ्गयोगेभ्यो
\item सरस्वतीसिद्धेभ्यो
\item लोकशिक्षकेभ्यो
\item अणिमादिसिद्धेभ्यो
\item यजमानेभ्यो
\item शापानुग्रहशक्तेभ्यो
\item जीवन्मुक्तेभ्यो
\item याजकेभ्यो
\item स्वतन्त्रशक्तिभ्यो
\item शिवपूजारतेभ्यो
\item ऋत्विग्भ्यो
\item स्वाधीनचित्तेभ्यो
\item जितेन्द्रियेभ्यो
\item शान्तेभ्यो
\item दान्तेभ्यो
\item शुद्धेभ्यो
\item मुनिमुख्येभ्यो
\item रुद्राक्षधारिभ्यो
\item सुमुख्येभ्यो
\item वल्कलाजिनधारिभ्यो
\item जटिलेभ्यो
\item तितिक्षुभ्यो
\item आस्तिकेभ्यो
\item कमण्डलुधारिभ्यो
\item उपरतेभ्यो
\item विप्रेभ्यो
\item सपत्नीकेभ्यो
\item श्रद्धालुभ्यो
\item साङ्गेभ्यो
\item द्विजेभ्यो
\item विष्णुभक्तेभ्यो
\item ब्राह्मणेभ्यो
\item वेदवेद्येभ्यो
\item विवेकिभ्यो
\item उपवीतिभ्यो
\item स्मृतिकर्तृभ्यो
\item विज्ञेभ्यो
\item मेधाविभ्यो
\item ब्रह्मिष्ठेभ्यो
\item पवित्रपाणिभ्यो
\item मन्त्रकृद्भ्यो
\item दीनबन्धुभ्यो
\item ब्रह्मनिष्ठेभ्यो
\item संस्कृतेभ्यो
\item भगवद्भ्यो
\item सत्कृतेभ्यो
\item भस्मधारिभ्यो
\item सुकृतिभ्यो
\item श्रीकश्यपादिसप्तमहर्षिभ्यो
\item अरुन्धत्यादिसप्तर्षिपत्नीभ्यो
\end{enumerate}
\end{multicols}
\textbf{...भ्यो नमः नानाविधपरिमलपत्रपुष्पाणि समर्पयामि।}

\twolineshloka*
{वनस्पतिरसोद्भूतो गन्धाढ्यस्सुमनोहरः}
{आत्रेयस्सर्वदेवानां धूपोऽयं प्रतिगृह्यताम्}
\textbf{...भ्यो नमः धूपमाघ्रापयामि।}

\twolineshloka*
{साज्यं त्रिवर्तिसंयुक्तं वह्निना योजितं मया}
{गृहाण मङ्गलं दीपं त्रैलोक्यतिमिरापहम्}
\textbf{...भ्यो नमः अलङ्कारदीपं सन्दर्शयामि।}

\ifbool{veda}{ॐ भूर्भुवः॒ सुवः॑। + ब्र॒ह्मणे॒ स्वाहा᳚।}{}
\twolineshloka*
{नानापक्वान्नसंयुक्तं रसैस्सप्तभिस्समन्वितम्}
{गृह्णन्तु मुनयस्सर्वे नैवेद्यं वसुदायकाः}
\textbf{...भ्यो नमः नैवेद्यं निवेदयामि।}
अमृतापिधानमसि। निवेदनानन्तरम् आचमनीयं समर्पयामि।

\twolineshloka*
{रसालजम्बूद्राक्षेक्षुकपित्थकदलीफलम्}
{खर्जूरपनसादीनि नालिकेरफलानि च}
\textbf{...भ्यो नमः नानाविधफलानि समर्पयामि। ताम्बूलं समर्पयामि।}

\twolineshloka*
{नीराजनं मुनिश्रेष्ठाः कोटिसूर्यसमप्रभाः}
{अहं भक्त्या प्रदास्यामि स्वीकुरुध्वम् ऋषीश्वराः}
\textbf{...भ्यो नमः कर्पूरनीराजनं दर्शयामि। नीराजनानन्तरम् आचमनीयं समर्पयामि।}

\twolineshloka*
{यानि कानि च पापानि ब्रह्महत्यासमानि च}
{तानि तानि विनश्यन्ति प्रदक्षिणपदे पदे}
\textbf{...भ्यो नमः प्रदक्षिणानि समर्पयामि।}

\twolineshloka*
{नमो महर्षिवृन्देभ्यो देवर्षिभ्यो नमो नमः}
{सर्वपापहरेभ्योऽस्तु वेदविद्भ्यो नमो नमः}
\textbf{...भ्यो नमः नमस्कारान् समर्पयामि।}

\twolineshloka*
{साक्षसूत्राश्च मुनयस्सर्वे वल्कलधारिणः}
{मन्त्रपुष्पं प्रदास्यामि स्वीकुरुध्वं मुनीश्वराः}
\textbf{...भ्यो नमः मन्त्रपुष्पं समर्पयामि। सुवर्णपुष्पं समर्पयामि।}

\sect{प्रार्थना}

\twolineshloka*
{एते सप्तर्षयस्सर्वे भक्त्या सम्पूजिता मया}
{सर्वे पापं व्यपोहन्तु ज्ञानतोऽज्ञानतः कृतम्}
\textbf{प्रार्थनाः समर्पयामि।}

\dnsub{अर्घ्यम्}
\resetShloka
ममोपात्त-समस्त-दुरित-क्षयद्वारा अरुन्धती-सहित-कश्यपादि-सप्तर्षि-प्रीत्यर्थं
ऋषिपूजान्ते क्षीरार्घ्यप्रदानम् उपायनदानं च करिष्ये॥

\twolineshloka
{प्रसन्नहृदयाश्शान्ता जटामण्डलमण्डिताः}
{भक्त्या चार्घ्यं मया दत्तं स्वीकुर्वन्तु मुनीश्वराः}
\textbf{अरुन्धती-सहित-कश्यपादि-सप्तर्षिभ्यो नमः इदमर्घ्यम्।} (त्रिः)

\twolineshloka
{जितकामाद्यरिचयाश्शमादिगुणभूषिताः}
{इदमर्घ्यं प्रदास्यामि प्रसन्नास्सन्तु मे सदा}
\textbf{इदमर्घ्यम्।} (त्रिः)

\twolineshloka
{ऋषयो ज्ञानसंसिद्धा ब्रह्मिष्ठा ब्रह्मवादिनः}
{इदमर्घ्यं प्रदास्यामि मत्पापशमनाय च}
\textbf{इदमर्घ्यम्।} (त्रिः)

\twolineshloka
{ब्रह्मदण्डधराश्शान्ताः कर्मठा वीतकल्मषाः}
{मत्पापपरिहारार्थम् इदमर्घ्यं मुनीश्वराः}
\textbf{इदमर्घ्यम्।} (त्रिः)

अनेन अर्घ्यप्रदानेन भगवन्तस्सर्वात्मकाः अरुन्धती-सहित-कश्यपादि-सप्तर्षयः प्रीयन्ताम्।

\dnsub{उपायन-दानम्}

\twolineshloka*
{सपात्रं फलसंयुक्तं सघृतं दक्षिणान्वितम्}
{उपायनं प्रदास्यामि व्रतसम्पूर्तिहेतवे}

\twolineshloka*
{भवन्तः प्रतिगृह्णन्तु तेजोरूपास्तपोधनाः}
{ऋषयस्तारकास्सन्तूपायनस्य प्रदानतः}

इति उपायनदानं कृत्वा।

\twolineshloka*
{न्यूनातिरिक्तकर्माणि मया यानि कृतानि च}
{क्षमध्वं तानि सर्वाणि यूयं सर्वे तपोधनाः}

इति ऋषिपूजाविधिः सम्पूर्णः॥

\sect{उद्यापन-विधिः}

एवं प्रतियामं (प्रतिवर्षं) पूजां कृत्वा उपायनानि दत्त्वा कथां च श्रुत्वा, सप्तमे
वर्षे ततः परेद्युः नित्यकर्म समाप्य ऋषिपञ्चमीव्रतोद्यापनहोमं करिष्य इति सङ्कल्प्य
स्वगृह्योक्तविधानेन अग्निमुखान्तं कृत्वा "सहस्रोमा" इति मन्त्रेण पायसान्नं
समिदाज्यैः अष्टोत्तरशतसंख्यया अष्टाविंशत्या वा जुहुयात्। ततः स्विष्टकृदादि-होमशेषं
समापयेत्। पुनः पूजां करिष्य इति सङ्कल्प्य आसनादि-षोडशोपचारपूजां कृत्वा,
"यज्ञेन यज्ञम्" इति मन्त्रेण ऋषीन् उद्वास्य, "तत्त्वायामि" इति वरुणम् उद्वास्य,
कुम्भजलेन "देवस्य त्वा" इति मन्त्रेण त्रिपद-गायत्र्या व्याहृतिभिश्च यजमानं
परिवारं च प्रोक्षयेत्, कुम्भजलेन यजमानम् अभिषिच्य, धौतवस्त्र-पुण्ड्रधारणानि कृत्वा
तीर्थं प्राश्य, कलश-वस्त्र-प्रतिमादानानि करिष्य इति सङ्कल्प्य,

\dnsub{दान-मन्त्रः}
\twolineshloka*
{प्रतिमां वस्त्रसंयुक्तां कुम्भोपकरणैर्युताम्}
{तुभ्यं दास्यामि विप्रेन्द्र यथोक्तफलदो भव}

इमां कश्यप-प्रतिमां सदक्षिणां सताम्बूलां कश्यप-प्रीतिं कामयमानः तुभ्यम् अहं
सम्प्रददे। एवं क्रमशः अत्रि-भरद्वाज-विश्वामित्र-गौतम-जमदग्नि-वसिष्ठ-अरुन्धती-
प्रतिमाः प्रीतिं कामयमानः दत्वा। अनन्तरं दशदान-षड्दान-मनोद्दिष्ट-दानानि
स्वशक्त्या विस्तरतः कृत्वा, दम्पति-पूजनं कृत्वा, ब्राह्मणान् भोजयित्वा,
ताम्बूलानन्तरम् आशिषः प्राप्य, तैः अनुज्ञातः पुत्रमित्रकलत्रादिभिः सह स्वयं
भुञ्जीत।

इति ब्रह्माण्डपुराणोक्त-ऋषिपञ्चमी-व्रतोद्यापन-विधिः सम्पूर्णः॥

\closesub

\sect{प्रार्थना}

\twolineshloka*
{यन्मयाऽऽचरितं देव व्रतमेतत् सुदुर्लभम्}
{गणेश त्वं प्रसन्नः सन् सफलं कुरु सर्वदा}

\sect{अपराध-क्षमापनम्}

\twolineshloka*
{यस्य स्मृत्या च नामोक्त्या तपः पूजाक्रियादिषु}
{न्यूनं सम्पूर्णतां याति सद्यो वन्दे गजाननम्}

\fourlineindentedshloka*
{कायेन वाचा मनसेन्द्रियैर्वा}
{बुद्‌ध्याऽऽत्मना वा प्रकृतेः स्वभावात्}
{करोमि यद्यत् सकलं परस्मै}
{नारायणायेति समर्पयामि}

\centerline{ॐ तत्सद्ब्रह्मार्पणमस्तु।}

\ifbool{katha}{\sect{ऋषिपञ्चमी-व्रत-कथा}
\centerline{\small{(मूलम्—श्रीभविष्योत्तरपुराणम्, यथा श्रीव्रतरत्नाकरे संगृहीतम्)}}

\twolineshloka
{श्रुतानि देवदेवेश व्रतानि सुबहूनि च}
{साम्प्रतं मे समाचक्ष्व व्रतं पापप्रणाशनम्} %॥१॥

\uvacha{ब्रह्मोवाच}

\twolineshloka
{शृणु राजन् प्रवक्ष्यामि व्रतानामुत्तमं व्रतम्}
{ऋषिपञ्चमीति विख्यातं सर्वपापहरं परम्} %॥२॥

\twolineshloka
{येन चीर्णेन राजेन्द्र नरकं नैव पश्यति}
{अत्रैवोदाहरन्तीममितिहासं पुरातनम्} %॥३॥

\twolineshloka
{वैदर्भे च द्विजवर उदको नाम नामतः}
{तस्य भार्या सुशीलेति पतिव्रतपरायणा} %॥४॥

\twolineshloka
{तस्या अपत्ययुगलं पुत्रो हि सुविभीषणः}
{अधीतवान् सुतस्तस्य वेदान् साङ्गपदक्रमान्} %॥५॥

\twolineshloka
{समानेन कुलीनेन सुता चापि विवाहिता}
{नवोढात्वेऽपि सा दैवाद्वैधव्यं प्राप सत्तमा} %॥६॥

\twolineshloka
{सतीत्वं पालयन्ती सा आस्ते निजपितुर्गृहे}
{तस्या दुःखेन सन्तप्तस्सुतं संस्थाप्य वेश्मनि} %॥७॥

\twolineshloka
{गङ्गातीरवनं प्राप्तः कदाचित्तु तया सह}
{स तत्राध्यापयामास शिष्यान् वेदं द्विजोत्तमः} %॥८॥

\twolineshloka
{सुता च कुरुते तस्य पितुश्शुश्रूषणं परम्}
{पितुश्शुश्रूषणं कृत्वा परिश्रान्ता कदाचन} %॥९॥

\twolineshloka
{निशीथे किल सुप्तायां क्रिमिराशिरजायत}
{शिष्यास्तथाविधां दृष्ट्वा विवस्त्रां प्रस्तरे स्थिताम्} %॥१०॥

\twolineshloka
{गत्वा निवेदयामासुस्तन्मातुः करुणां गताः}
{न जानीमो वयं किञ्चिद्देवि साध्वी तथाविधा} %॥११॥

\twolineshloka
{क्रिमिराशिमयी जाता मातस्सम्प्रति दृश्यते}
{वज्रपातसमं श्रुत्वा तच्छिष्यैस्समुदाहृतम्} %॥१२॥

\twolineshloka
{सा श्रान्तमानसा शीघ्रं तत्समीपमुपागमत्}
{सुतां तथाविधां दृष्ट्वा विललाप सुदुःखिता} %॥१३॥

\twolineshloka
{उरश्च ताडयामास सुतरां मोहमाप च}
{क्षणेन प्राप्य चैतन्यमुत्थाप्य परिमृज्य च} %॥१४॥

\twolineshloka
{समालिङ्ग्य च बाहुभ्यां निन्ये तत्पितुरन्तिकम्}
{स्वामिन् कथय मे साध्वी केन दुष्कृतकर्मणा} %॥१५॥

\twolineshloka
{निशीथे च प्रसुप्तेयं जायते क्रिमिसङ्कुला}
{एवं श्रुत्वा ततो वाक्यम् ऋषिर्ध्यानपरायणः} %॥१६॥

\onelineshloka
{ज्ञात्वा निवेदयामास तस्याः प्राग्जन्मचेष्टितम्} %॥१७॥

\uvacha{ऋषिरुवाच}

\twolineshloka
{प्राणीयं सप्तमेऽतीते जन्मनि ब्राह्मणी ह्यभूत्}
{रजस्वला च सञ्जाता भाण्डादीनस्पृशत्तदा} %॥१८॥

\twolineshloka
{अस्यास्तु पाप्मना तेन जायते क्रिमिवद्वपुः}
{रजस्वलायाः पापेन युक्ता भवति साऽनघे} %॥१९॥

\twolineshloka
{प्रथमेऽहनि चण्डाली द्वितीये ब्रह्मघातिनी}
{तृतीये रजकी प्रोक्ता चतुर्थेऽहनि शुद्ध्यति} %॥२०॥

\twolineshloka
{तदा तया सखीसङ्घातं दृष्ट्वाऽवमानितम्}
{दृष्टव्रतप्रभावेन जाता द्विजकुलेऽमले} %॥२१॥

\twolineshloka
{अवमानाद्व्रतस्यास्य क्रिमिराशिरिवाधुना}
{एतत्ते कथितं सर्वं कारणं दुष्कृतस्य च} %॥२२॥

\uvacha{सुशीलोवाच}

\twolineshloka
{दर्शनादपि जन्तोरस्य द्विजानां निर्मले कुले}
{जन्म युष्मद्विधानां हि जायते ब्रह्मतेजसाम्} %॥२३॥

\twolineshloka
{अवज्ञया प्रजायन्ते निशीथे क्रिमिराशयः}
{महाश्चर्यं कथं नाथ तद्व्रतं कथयस्व मे} %॥२४॥

\uvacha{उदक उवाच}

\twolineshloka
{सुशीले शृणु तत्सम्यक् व्रतं नामोत्तमं व्रतम्}
{येन चीर्णेन सहसा पापादस्मात् प्रमुच्यते} %॥२५॥

\twolineshloka
{दुःखत्रयाच्च मुच्येत नारी सौभाग्यमाप्नुयात्}
{कल्याणानि विवर्धन्ते सम्पदश्च निरापदः} %॥२६॥

\uvacha{युधिष्ठिर उवाच}

\twolineshloka
{श्रुतानि देवदेवेश व्रतानि सुबहूनि च}
{साम्प्रतं त्वन्यदाचक्ष्व व्रतं पापप्रणाशनम्} %॥२७॥

\uvacha{श्रीकृष्ण उवाच}

\twolineshloka
{अथान्यदपि राजेन्द्र पञ्चम्यामृषिपूजनम्}
{कथयिष्यामि यत्कृत्वा नारी पापात् प्रमुच्यते} %॥२८॥

\uvacha{युधिष्ठिर उवाच}

\twolineshloka
{कीदृशी पञ्चमी कृष्ण कथं तदृषिपूजनम्}
{पातकान्मुच्यते कस्मान्नारी यदुकुलोद्भव} %॥२९॥

\twolineshloka
{पापानि च बहून्यत्र विद्यन्ते किल केशव}
{कथं वा ऋषिपञ्चम्यां नारी कस्मात् प्रमुच्यते} %॥३०॥

\uvacha{श्रीकृष्ण उवाच}

\twolineshloka
{अज्ञानाज्ज्ञानतो वापि या स्त्री जाता रजस्वला}
{दुष्टा स्पृशति भाण्डानि गृहकर्मणि संस्थिता} %॥३१॥

\twolineshloka
{प्राप्नोति सा महापापं सर्वा ना रजस्वलाः}
{शृणु तत्कारणं येन वर्जनीया रजस्वला} %॥३२॥

\twolineshloka
{प्रोत्सार्या गृहतो दूरं चातुर्वर्ण्येन भारत}
{ब्रह्महत्यां पुरा शक्रो वृत्रं हत्वा ह्यवाप च} %॥३३॥

\twolineshloka
{तया वै राजशार्दूल पीडितो वृत्रसूदनः}
{ब्रह्माणं समुपागच्छदात्मशुद्धिकारणात्} %॥३४॥

\twolineshloka
{शुद्ध्यै शक्रस्य राजेन्द्र प्रहृष्टेनान्तरात्मना}
{विभज्य ब्रह्महत्यां तु चतुर्धा स चतुर्मुखः} %॥३५॥

\twolineshloka
{प्राक्षिपद्राजशार्दूल चतुःस्थानेषु वै तदा}
{वह्नौ प्रथमजिह्वासु नदीनां प्रथमोदके} %॥३६॥

\twolineshloka
{गिरिवृक्षादिभूमौ च नारीरजसि पार्थिव}
{ततो रजस्वला नारी प्रोत्सार्या च प्रयत्नतः} %॥३७॥

\twolineshloka
{ब्राह्मणशासनात्पार्थ चातुर्वर्ण्येन सर्वदा}
{प्रथमेऽहनि चण्डाली द्वितीये ब्रह्मघातिनी} %॥३८॥

\twolineshloka
{तृतीये रजकी प्रोक्ता चतुर्थेऽहनि शुद्ध्यति}
{अज्ञानाज्ज्ञानतो वापि जातं सम्पर्कपातकम्} %॥३९॥

\twolineshloka
{तत्पापसंक्षयार्थं वै कार्येयमृषिपञ्चमी}
{सर्वपापप्रशमनी सर्वोपद्रवनाशिनी} %॥४०॥

\onelineshloka
{ब्रह्मक्षत्रियविट्शूद्रस्त्रीभिः कार्या विशेषतः} %॥४१॥

\uvacha{कृष्ण उवाच}

\onelineshloka
{अत्रार्थे यत्पुरा वृत्तं प्रवक्ष्यामि कथात्मकम्} %॥४२॥

\twolineshloka
{पुरा कृतयुगे राजा विदर्भाणां बभूव ह}
{श्येनजिन्नाम राजर्षिश्चातुर्वर्ण्यानुपालकः} %॥४३॥

\twolineshloka
{तस्य देशे वसन्विप्रो वेदवेदाङ्गपारगः}
{सुमित्रो नाम विप्रेन्द्रः सर्वभूतहिते रतः} %॥४४॥

\twolineshloka
{कृषिवृत्त्या सदा युक्तः कुटुम्बपरिपालकः}
{तस्य भार्या सुसाध्वी च पतिशुश्रूषणे रता} %॥४५॥

\twolineshloka
{जयश्रीर्नाम विख्याता बहुभृत्यसुहृज्जना}
{अतिचिन्तान्विता सा च प्रावृट्काले सुमध्यमा} %॥४६॥

\twolineshloka
{क्षेत्रादिषु रता साध्वी चाकुलीकृतमानसा}
{एकदा साऽऽत्मनः प्राप्तमृतुकालं व्यलोकयत्} %॥४७॥

\twolineshloka
{रजस्वलाऽपि सा राजन् गृहकर्म चकार ह}
{भाण्डादीन्यस्पृशद्राजन् प्राप्तर्तुरपि भामिनी} %॥४८॥

\twolineshloka
{कालेन बहुना साध्वी पञ्चत्वमगमत्तदा}
{तस्या भर्ताऽपि विप्रोऽसौ कालधर्ममुपेयिवान्} %॥४९॥

\twolineshloka
{एवं तौ दम्पती राजन् स्वकर्मवशगौ तदा}
{भार्या तस्य जयश्रीस्सा ऋतुसम्पर्कदोषतः} %॥५०॥

\twolineshloka
{शुनीयोनिसमुत्पन्ना शुनी चाभून्नरेश्वर}
{तस्यास्सम्पर्कदोषेण बलीवर्दो बभूव सः} %॥५१॥

\twolineshloka
{ऋतुसम्पर्कदोषेण तिर्यग्योनिमुपागतौ}
{स्वधर्माचरणाज्जातावुभौ जातिस्मरौ तदा} %॥५२॥

\twolineshloka
{सुतस्यैव गृहे राजन् स्मरन्तौ पूर्वपातकम्}
{सुमित्रस्य च पुत्रोऽभूद्गुरुशुश्रूषणे रतः} %॥५३॥

\twolineshloka
{सुमतिर्नाम धर्मज्ञो देवतातिथिपूजकः}
{कालक्षयाहे सम्प्राप्ते पितुस्तु सुमतिस्तदा} %॥५४॥

\twolineshloka
{भार्यां चन्द्रवतीं प्राप्य पैतृकश्रद्धयाऽन्वितः}
{अद्य साम्वत्सरदिनं पितुर्मे चारुहासिनि} %॥५५॥

\twolineshloka
{भोजनीया द्विजा भीरु पाकशुद्धिर्विधीयताम्}
{तया कृता पाकशुद्धिस्सुमते शासनाद्गृहे} %॥५६॥

\twolineshloka
{मुक्तं पायसभाण्डे वै सर्पेण गरलं ततः}
{दृष्ट्वा ब्रह्मवधाद्भीता शुनी भाण्डानि साऽस्पृशत्} %॥५७॥

\twolineshloka
{द्विजभार्या च तां दृष्ट्वा उल्मुकेन जघान ह}
{भाण्डादीनि च प्रक्षाल्य त्यक्त्वा पाकं सुमध्यमा} %॥५८॥

\twolineshloka
{पुनः पाकं च कृत्वा तु श्राद्धं कृत्वा विधानतः}
{ततो भुक्तेषु विप्रेषु नोच्छिष्टं च ददौ बहिः} %॥५९॥

\twolineshloka
{भूमौ क्षिप्तं तया शुन्या उपवासस्ततोऽभवत्}
{ततो रात्र्यां प्रवृत्तायां सा शुनी क्षुधिता भृशम्} %॥६०॥

\onelineshloka
{बलीवर्दमुपागत्य भर्तारमिदमब्रवीत्} %॥६१॥

\uvacha{शुन्युवाच}

\twolineshloka
{बुभुक्षिता गृहे भर्तर्न दत्तं भोजनादिकम्}
{ग्रासादिकं च न प्राप्तं क्षुधा मां बाधते भृशम्} %॥६२॥

\twolineshloka
{अन्यस्मिन् दिवसे पुत्रो मम लेह्यं ददात्यसौ}
{अद्य मह्यं किमप्येष उच्छिष्टमपि नो ददौ} %॥६३॥

\twolineshloka
{पायसान्ने पपाताद्य गरलं सर्पसम्भवम्}
{मया विचिन्त्य मनसा मरिष्यन्ति द्विजोत्तमाः} %॥६४॥

\twolineshloka
{संस्पृष्टं पायसं गत्वा तत्राहं ताडिता भृशम्}
{दुःखितं तेन मे गात्रं कटिर्भग्ना करोमि किम्} %॥६५॥

\uvacha{बलीवर्द उवाच}

\onelineshloka
{भद्रे ते पापजं फलम्} %॥६६॥

\twolineshloka
{किं करोमि शक्तोऽहं भारवाहोऽद्य सुस्थितः}
{अद्याहमात्मजक्षेत्रे वाहितस्सकलं दिनम्} %॥६७॥

\twolineshloka
{मारितश्चात्मजेनाहं मुखं बद्ध्वा बुभुक्षितः}
{वृथा श्राद्धं कृतं तेन जाताऽद्य मम कष्टता} %॥६८॥

\uvacha{श्रीकृष्ण उवाच}

\twolineshloka
{तयोस्संवदतोरेव मातापित्रोश्च भारत}
{श्रुत्वा पुत्रस्तदा वाक्यं यदुक्तं च तदोभयोः} %॥६९॥

\twolineshloka
{ततो रजन्यां तत्कालं ददौ तस्यैव भोजनम्}
{पितरौ तौ विदित्वा तु दत्तवान् सुमतिस्तदा} %॥७०॥

\twolineshloka
{ततोऽसौ दुःखितः पुत्रो ज्ञात्वाऽवस्थां तदा तयोः}
{मातापित्रोस्तु राजेन्द्र दुःखात्मा प्रस्थितो वनम्} %॥७१॥

\twolineshloka
{ज्ञातुमिच्छामि कष्टं वै मातापित्रोश्च भारत}
{तत्र गत्वा ज्ञानवृद्धानृषीन् सप्त वनाश्रितान्} %॥७२॥

\onelineshloka
{प्रणिपत्याब्रवीद्वाक्यं हितं चैव तदा तयोः} %॥७३॥

\uvacha{सुमतिरुवाच}

\twolineshloka
{कथयध्वं विप्रवर्याः प्रश्नमेकं समाहिताः}
{केन कर्मविपाकेन पितरौ मे तपोधनाः} %॥७४॥

\onelineshloka
{इमामवस्थां सम्प्राप्तौ मुच्येते पातकात्कथम्} %॥७५॥

\uvacha{श्रीकृष्ण उवाच}

\twolineshloka
{तदाकर्ण्य वचस्तस्य सुमतेर्दुःखितस्य च}
{ऋषिस्सर्वतपा नाम सर्वज्ञः करुणान्वितः} %॥७६॥

\onelineshloka
{सुमतिं प्रत्युवाचेदं तत्पित्रोर्मुक्तये तदा} %॥७७॥

\uvacha{ऋषिरुवाच}

\onelineshloka
{तव माता पुरा विप्र स्वगृहे बालभावतः} %॥७८॥

\twolineshloka
{ऋतुं प्राप्तं विदित्वा तु सम्पर्कमकरोद्द्विज}
{तेन कर्मविपाकेन शुनीयोनिमुपागता} %॥७९॥

\twolineshloka
{पिताऽपि स्पर्शदोषेण बलीवर्दो बभूव ह}
{एतयोर्मुक्तिकामार्थं कुरु त्वमृषिपञ्चमीम्} %॥८०॥

\twolineshloka
{भार्यया सह विप्रेन्द्र ऋषीन् सम्पूज्य यत्नतः}
{आचरस्व व्रतं तत्र सप्तवर्षं द्विजोत्तम} %॥८१॥

\twolineshloka
{अन्ते चोद्यापनं कुर्याद्वित्तशाठ्यविवर्जितः}
{शाकाहारस्तु कर्तव्यो नीवारैश्श्यामकैस्तदा} %॥८२॥

\twolineshloka
{कन्दं वाऽथ फलं मूलं हलकृष्टं न भक्षयेत्}
{प्राप्य भाद्रपदे मासि शुक्लपक्षस्य पञ्चमीम्} %॥८३॥

\twolineshloka
{तस्यां मध्याह्नसमये नद्यादौ विमले जले}
{कृत्वाऽपामार्गसमिधा दन्तधावनमादितः} %॥८४॥

\twolineshloka
{आयुर्बलं यशो वर्चः प्रजाः पशुवसूनि च}
{ब्रह्म प्रज्ञां च मेधां च त्वन्नो देहि वनस्पते} %॥८५॥

\twolineshloka
{संप्रार्थ्यानेन मन्त्रेण कुर्याद्वै दन्तधावनम्}
{मुखदुर्गन्धनाशाय दन्तानां च विशुद्धये} %॥८६॥

\twolineshloka
{ष्ठीवनाय च गात्राणां कुर्वेऽहं दन्तधावनम्}
{अनेन दन्तान् संशोध्य स्नायान्मृत्स्नानपूर्वकम्} %॥८७॥

\twolineshloka
{तिलामलककल्केन केशान् संशोध्य यत्नतः}
{परिधाय नवं शुद्धं वासोयुग्मं समाहितः} %॥८८॥

\twolineshloka
{विधाय नित्यकर्माणि गृहं गत्वा तान् ऋषीन्}
{स्थापयेद्विधिवद्भक्त्या पञ्चामृतरसैश्शुभैः} %॥८९॥

\twolineshloka
{चन्दनागरुकर्पूरैर्विलिप्य च सुगन्धिभिः}
{पूजयेद्विविधैः पुष्पैर्गन्धधूपादिदीपकैः} %॥९०॥

\twolineshloka
{समाच्छाद्य शुभैर्वस्त्रैस्सोपवीतैर्यथाविधि}
{ततो नैवेद्यसम्पन्नैर्दद्याच्छुभैः फलैः} %॥९१॥

\onelineshloka
{पूजयस्व ऋषीन् दिव्यान् अरुन्धत्या समन्वितान्} %॥९२॥

\twolineshloka
{कश्यपोऽत्रिर्भरद्वाजो विश्वामित्रोऽथ गौतमः}
{जमदग्निर्वसिष्ठश्च साध्वी चैवाप्यरुन्धती} %॥९३॥

\twolineshloka
{मन्त्रेणानेन सप्तर्षीन् पूजयेत्सुसमाहितः}
{व्रतेन ऋषिपञ्चम्याः कृतेनैव द्विजोत्तम} %॥९४॥

\onelineshloka
{ऋतुसम्पर्कजो दोषः क्षयं याति न संशयः} %॥९५॥

\uvacha{श्रीकृष्ण उवाच}

\onelineshloka
{तच्छ्रुत्वा सुमतिर्वाक्यं परमं मुनिभाषितम्} %॥९६॥

\twolineshloka
{गृहमेत्य व्रतं चक्रे सभार्यः श्रद्धयाऽन्वितः}
{व्रतं तु ऋषिपञ्चम्यास्सर्वपापप्रणाशनम्} %॥९७॥

\twolineshloka
{कृत्वा सर्वं यथोक्तं च मातापित्रोः फलं ददौ}
{व्रतपुण्यप्रभावेन माता तस्य श्वयोनितः} %॥९८॥

\twolineshloka
{मुक्ता नृपतिशार्दूल विमानवरमास्थिता}
{दिव्याम्बरधरा भूत्वा गता स्वर्गं च भारत} %॥९९॥

\twolineshloka
{पिताऽपि स मृतोऽगच्छत् सुमतेः पशुयोनितः}
{स्वर्गं प्राप्तो महाराज व्रतस्यास्य प्रभावतः} %॥१००॥

\twolineshloka
{दिव्याम्बरधरो भूत्वा गतस्स्वर्गं स भारत}
{कायिकं वाचिकं वापि मानसं यच्च दुष्कृतम्} %॥१०१॥

\twolineshloka
{तत्सर्वं विलयं याति व्रतस्यास्य प्रभावतः}
{तस्य यज्जायते पुण्यं तच्छृणुष्व नृपोत्तम} %॥१०२॥

\twolineshloka
{सर्वव्रतेषु यत्पुण्यं सर्वतीर्थेषु यत्फलम्}
{सर्वदानेषु दत्तेषु तदेतद्व्रतचारणात्} %॥१०३॥

\twolineshloka
{कुरुते या व्रतं नारी सा भवेत्सुखभागिनी}
{रूपलावण्ययुक्ता च पुत्रपौत्राभिसंयुता} %॥१०४॥

\twolineshloka
{इह लोके सुखं यायात् परत्र च परां गतिम्}
{एतत्ते कथितं राजन् व्रतानामुत्तमं व्रतम्} %॥१०५॥

\twolineshloka
{सर्वसम्पत्प्रदं चैव नारीणां पापनाशनम्}
{धनं यशश्च स्वर्गं च पुत्रानपि युधिष्ठिर} %॥१०६॥

\onelineshloka
{पठतां शृण्वतां चापि दद्यात्पापप्रणाशनम्} %॥१०७॥

॥ इति श्रीभविष्योत्तरपुराणे ऋषिपञ्चमीव्रतकथा सम्पूर्णा ॥
}{}

\closesection
